% Options for packages loaded elsewhere
\PassOptionsToPackage{unicode}{hyperref}
\PassOptionsToPackage{hyphens}{url}
\documentclass{article}
\usepackage[a4paper,margin=1in]{geometry}
\usepackage{xcolor}
\usepackage{amsmath,amssymb}
\setcounter{secnumdepth}{3}
\usepackage{iftex}
\ifPDFTeX
  \usepackage[T1]{fontenc}
  \usepackage[utf8]{inputenc}
  \usepackage{textcomp} % provide euro and other symbols
\else % if luatex or xetex
  \usepackage{unicode-math} % this also loads fontspec
  \defaultfontfeatures{Scale=MatchLowercase}
  \defaultfontfeatures[\rmfamily]{Ligatures=TeX,Scale=1}
\fi
\usepackage{lmodern}
\ifPDFTeX\else
  % xetex/luatex font selection
\fi
% Use upquote if available, for straight quotes in verbatim environments
\IfFileExists{upquote.sty}{\usepackage{upquote}}{}
\IfFileExists{microtype.sty}{% use microtype if available
  \usepackage[]{microtype}
  \UseMicrotypeSet[protrusion]{basicmath} % disable protrusion for tt fonts
}{}
\makeatletter
\@ifundefined{KOMAClassName}{% if non-KOMA class
  \IfFileExists{parskip.sty}{%
    \usepackage{parskip}
  }{% else
    \setlength{\parindent}{0pt}
    \setlength{\parskip}{6pt plus 2pt minus 1pt}}
}{% if KOMA class
  \KOMAoptions{parskip=half}}
\makeatother
\setlength{\emergencystretch}{3em} % prevent overfull lines
\providecommand{\tightlist}{%
  \setlength{\itemsep}{0pt}\setlength{\parskip}{0pt}}
\usepackage{bookmark}
\usepackage{csquotes}
\usepackage[backend=biber,style=authoryear]{biblatex}
\addbibresource{references.bib}
\IfFileExists{xurl.sty}{\usepackage{xurl}}{} % add URL line breaks if available
\urlstyle{same}
\hypersetup{
  pdftitle={Interim Report: Experience-Dependent Procedural Knowledge in GUI Agents: A Taxonomy and Research Agenda},
  hidelinks,
  pdfcreator={LaTeX via pandoc}}

\title{Interim Report\\[0.4em]\large Experience-Dependent Procedural Knowledge in GUI Agents: A Taxonomy and Research Agenda}
\author{Li Mengxiao}
\date{\today}

\begin{document}
\maketitle
\tableofcontents
\newpage

\section{Abstract}\label{1-abstract}

LLM-based GUI agents have recently achieved rapid gains in perception,
planning, and action execution across mobile, desktop, and web
environments. Yet much of this progress still relies on stronger
backbones, larger training corpora, longer context windows, or
app-specific workflow hints. This survey argues that the next bottleneck
is not generic GUI skill, but the lack of a reusable and revisable
experience layer: \textbf{experience-dependent procedural knowledge},
meaning action-relevant know-how acquired from actually interacting with
interfaces and reused later rather than within the same run only. We
study this problem through a taxonomy that asks \textbf{what} is
learned, \textbf{when} experience is written back, and \textbf{how}
experience is matched and applied. The taxonomy further separates
\texttt{Match\ Operator} from \texttt{Application\ Carrier}, that is,
how a system finds a past experience and how that experience actually
influences behavior. This distinction allows retrieval to be separated
from enforcement and enables current systems to be analyzed at the level
of capability patterns rather than whole-paper labels. Under this view,
existing GUI agents already instantiate parts of the \texttt{post-task}
memory space, especially
\texttt{post-task\ /\ Skills\ /\ similarity\ retrieval\ ×\ context\ injection},
but still rarely support rule-level generalization, cross-task
procedural reuse, or failure-driven memory rewrite. We then synthesize
evidence across GUI-agent, agent-memory, and self-evolving-agent
literatures to argue that two gaps now form the most defensible research
focus: procedural memory for GUI and failure-driven write-back. Episodic
memory is positioned as supporting infrastructure, while user-centric
memory remains a confirmed but secondary research direction. The
resulting survey contributes not only a literature organization scheme,
but also a methodological agenda: memory units should be structured
procedural rules rather than raw trajectories, retrieval and application
should be evaluated separately, and write-back should be treated as a
first-class design choice rather than an implementation detail. This
framing narrows the field\textquotesingle s broad memory rhetoric into a
more testable target for future GUI-agent research.

\begin{center}\rule{0.5\linewidth}{0.5pt}\end{center}

\section{Introduction}\label{2-introduction}

\subsection{Why This Problem
Matters}\label{21-why-this-problem-matters}

Graphical user interfaces remain the dominant interface layer for
real-world digital work. A practical agent that can reliably operate
software, websites, or mobile applications must therefore reason over
GUI state, choose grounded actions, and recover from interface-specific
failures. Recent LLM- and VLM-based systems have made clear progress on
all three fronts. They can parse screenshots or accessibility trees,
decompose tasks into sub-steps, and even exploit auxiliary memory to
improve efficiency or robustness. Yet the strongest current systems
still exhibit a recurring weakness: once the interface changes, the
workflow becomes unfamiliar, or a local mistake reveals a hidden
precondition, much of the useful knowledge has to be re-discovered
rather than reused.

This weakness is easy to underestimate because it is often masked by
gains from stronger backbones, larger context windows, or better
training pipelines. But these improvements do not by themselves answer a
more structural question: what should the agent remember after acting?
Some GUI knowledge is already likely encoded in the base model, such as
recognizing a button, a search box, or a common navigation pattern. What
remains unstable is the interaction-derived part: knowledge about which
submenu usually hides the needed function, which form field must be
completed before another becomes active, which recovery action avoids
losing state, or which failure pattern indicates that a previously
successful tactic no longer applies. This is not generic GUI
familiarity. It is experience-dependent procedural knowledge.

The literature reviewed in later sections suggests that this is now the
most useful way to state the bottleneck. Existing systems such as
AppAgent~\parencite{zhang2023appagent},
MobileGPT~\parencite{lee2024mobilegpt},
MAGNET~\parencite{sun2026magnet},
ActionEngine~\parencite{zhong2026actionengine},
IntentCUA~\parencite{lee2026intentcua}, and GUI-Owl-1.5 /
MobileAgentV3.5~\parencite{xu2026mobileagentv35} have shown that
app-specific knowledge documents, sub-task graphs, workflow memory,
structural state graphs, and intent abstractions can all help. They
have not yet shown that post-task experience can be consolidated into
cross-task reusable rules and revised after failure. The central
motivation of this survey is therefore not to ask whether GUI agents
have any memory, but whether they can build the right kind of memory: a
reusable and revisable experience layer that goes beyond static hints
and transient context management.

\subsection{Scope and Boundary}\label{22-scope-and-boundary}

This survey focuses on LLM-based GUI agents that externalize or
operationalize experience reuse. Concretely, it is concerned with memory
objects and write-back mechanisms that persist beyond a single run and
can influence later task execution. The main emphasis is on the
transition from \texttt{post-task} experience to \texttt{cross-task}
reuse, in other words from learning after one task finishes to using
that learning in later tasks. The survey therefore prioritizes systems
that either already instantiate some form of externalized experience
reuse in GUI environments, or provide transferable precedents for how
such reuse could be built.

The survey is not mainly about generic prompt engineering, pure
training-stage self-evolution, or one-shot improvements that never
survive beyond the current episode. Nor is it a general survey of
text-only agent memory. Work from adjacent non-GUI domains is included
only when it provides a precedent that directly informs the missing GUI
capability, such as workflow induction, long-term storage management,
experience abstraction, or memory rewrite. This boundary is important
because otherwise the discussion quickly collapses into an overly broad
claim that "memory helps agents." That statement is true but not
discriminative enough to guide methodology.

The key boundary throughout the paper is the following: we distinguish
reusable experience-dependent procedural knowledge from generic GUI
priors already likely encoded in the base model. This distinction
excludes many tempting but unhelpful examples. A model remembering that
search bars often appear near the top of a page is not yet evidence of
learned experience. A model learning that, in a family of settings
pages, the relevant option is often hidden behind a secondary overflow
path and that a failed save usually reflects a missing prerequisite
field is much closer to the target object of this survey.

\subsection{Research Question}\label{23-research-question}

\begin{quote}
\textbf{RQ}: How can LLM-based GUI agents represent interaction-derived,
experience-dependent procedural knowledge as retrievable, revisable, and
cross-task reusable memory, and improve it through failure-aware
write-back from post-task experience?
\end{quote}

This research question is narrower than a generic "memory for GUI
agents" framing and commits the survey to three specific requirements.
First, the memory must represent interaction-derived knowledge rather
than merely restating generic interface priors. Second, the memory must
be retrievable and reusable across future tasks rather than only within
the current trajectory. Third, the memory must be revisable, which means
failure is treated not just as a runtime event but as a signal for
memory maintenance. Put simply, the goal is not just to remember what
happened before, but to remember it in a form that can be reused and
corrected later. These three requirements separate the framing adopted
here from neighboring but insufficiently aligned lines such as raw
trajectory replay, app-specific workflow caching, or stronger
training-side data flywheels.

\subsection{Contributions of This
Survey}\label{24-contributions-of-this-survey}

This survey makes four concrete contributions. First, it proposes an
operational taxonomy over \texttt{What} is learned, \texttt{When}
experience is written back, and \texttt{How} experience is matched and
applied, instead of collapsing GUI skill, memory, and self-improvement
into one undifferentiated category. Second, it uses taxonomy coverage
together with counter-evidence analysis to separate valuable blanks from
merely unoccupied cells. Third, it argues that procedural memory for GUI
and failure-driven write-back form the most defensible research focus,
while episodic memory and user-centric personalization are better
treated as supporting and secondary directions, respectively. Fourth, it
translates these conclusions into a methodological agenda, so that the
survey can serve not only as a related-work chapter but also as a bridge
into subsequent method design.

\begin{center}\rule{0.5\linewidth}{0.5pt}\end{center}

\section{Background and Problem
Formulation}\label{3-background-and-problem-formulation}

\subsection{Three Source Threads}\label{31-three-source-threads}

The problem studied in this survey sits at the intersection of three
research threads that are often discussed separately. GUI-agent work
contributes the grounded task setting, memory research contributes the
storage and retrieval vocabulary, and self-evolving-agent work
contributes the mechanisms for turning experience into future
improvement. The central research focus only becomes clear when these
three threads are read together.

\subsubsection{GUI Agents}\label{311-gui-agents}

The GUI-agent literature establishes the practical environment in which
the problem matters. These systems must perceive a screen, infer the
relevant task state, and choose actions that remain grounded under
interface variation, multi-step dependencies, and partial observability.
The strongest recent systems show that meaningful progress is possible
through app-specific knowledge documents, hierarchical sub-task memory,
workflow memory, structural state graphs, and intent-level planning
abstractions, as illustrated by AppAgent~\parencite{zhang2023appagent},
MobileGPT~\parencite{lee2024mobilegpt},
MobileAgent-v2~\parencite{wang2024mobileagentv2},
MAGNET~\parencite{sun2026magnet},
ActionEngine~\parencite{zhong2026actionengine}, and
IntentCUA~\parencite{lee2026intentcua}. This body of work therefore
proves that GUI agents do not fail only because of weak perception.
They also fail because the environment reveals interaction-specific
regularities that are hard to retain and reuse.

At the same time, GUI-native work also exposes the current limits of
memory design. Many systems still store either app-bound semantic
artifacts, transient working summaries, workflow templates, or
structural priors. These are all useful, but they answer different
questions. Some help an agent become familiar with a single app. Others
help it finish the current run more reliably. Fewer address whether
experience can survive as a revisable cross-task asset. This is why the
GUI thread provides the task pressure and the failure modes, but not yet
a complete answer to the current research question.

\subsubsection{Agent Memory}\label{312-agent-memory}

The agent-memory literature contributes a more precise language for
talking about what is being stored and why. It shows that memory is not
a single object but a family of functions: episodic memory preserves
past events, semantic memory stores relatively stable world knowledge,
procedural memory supports action selection, and user-centric memory
tracks preference or personalization signals. It also provides concrete
mechanisms such as selective retrieval, reflection-triggered
abstraction, long-term storage management, and explicit memory
evolution, as seen in Generative
Agents~\parencite{park2023generativeagents},
MemGPT~\parencite{packer2023memgpt}, and
A-MEM~\parencite{xu2025amem}. Without this
literature, GUI memory discussions tend to stay at the level of "the
agent remembers something." With it, the survey can ask sharper
questions about whether a given GUI artifact is serving as working
memory, semantic documentation, procedural reuse, or only a loose prompt
aid.

This thread also matters because it removes a recurring conceptual
objection. Long-term memory is not just a larger prompt. Text-domain
systems such as Generative
Agents~\parencite{park2023generativeagents},
MemGPT~\parencite{packer2023memgpt}, and
A-MEM~\parencite{xu2025amem} show that memory can be externally stored,
selectively retrieved, reorganized, and even rewritten in response to
new evidence. None of these methods directly solves GUI interaction,
because their memory substrates are mostly textual and their actions
are not visually grounded. But they establish that retrieval,
reflection, storage management, and write-back are not speculative
ideas. They are already viable memory primitives that GUI-agent
research can adapt.

\subsubsection{Self-Evolving
Agents}\label{313-self-evolving-agents}

The self-evolving-agent literature contributes the missing dynamic
piece: how experience is turned into improved future behavior. This
thread is especially important because it shifts the question from
storage to transformation. AWM~\parencite{wang2024awm} shows that
trajectories can be distilled into reusable workflows instead of being
replayed verbatim. ExpeL~\parencite{zhao2024expel} shows how
success-failure contrasts can maintain an editable pool of insights.
SkillRL~\parencite{xia2026skillrl} shows that failures can be distilled
into generalized skills rather than left as raw cautionary traces.
AgentKB~\parencite{tang2025agentkb} shows that experience can transfer
across heterogeneous agent frameworks when retrieval is filtered and
interference is controlled. These works collectively demonstrate that
"learning from experience" can be implemented as a memory pipeline
without continual backbone retraining.

In this review, this thread matters less as direct GUI evidence than as
methodological precedent. It shows that workflow induction, experience
abstraction, and failure-aware refinement are feasible in neighboring
environments. The unresolved issue is therefore not whether such
mechanisms can exist at all, but whether they can be grounded in
multimodal GUI settings where failure often depends on visual ambiguity,
hidden state, or app-specific control logic. This is the bridge between
self-evolving agents and the present focus on procedural memory and
failure-driven write-back.

\subsection{Key Definitions}\label{32-key-definitions}

To keep the rest of the survey precise, four terms need to be used in an
explicitly operational sense.

\textbf{Experience-dependent procedural knowledge} refers to
action-guiding knowledge that is revealed through interaction and is not
adequately explained by generic model prior alone. It is not the fact
that buttons are clickable or that a save action often appears near the
top of a page. It is the knowledge that, under a certain interface
pattern, the effective route is usually hidden behind a secondary menu,
that a field activation depends on an unstated prerequisite, or that a
certain recovery action prevents state loss after a failed submission.

\textbf{Memory unit} refers to the smallest experience artifact worth
retrieving and revising. In the framing adopted here, a memory unit is
not a raw action trace, but a structured procedural rule with trigger,
procedure, constraint, failure signal, revision cue, and scope. In
simpler terms, it is the smallest reusable lesson the system can store.
This definition matters because it excludes two common but insufficient
extremes: verbatim trajectory replay on one side and over-general prompt
advice on the other.

\textbf{Write-back} refers to the act of consolidating post-task
evidence into the long-term memory store so that future behavior can
change. Write-back may add a new memory unit, edit an existing one,
split an over-broad rule, downweight a memory whose scope was
overestimated, or rewrite a previously stored policy. The key point is
that write-back is about memory maintenance after task outcome is known,
not merely about using memory during inference.

\textbf{Cross-task reuse} refers to any reuse setting in which memory
created from one task instance influences another task instance beyond
literal replay. The most important levels here are repeated-task reuse,
same-app cross-task reuse, and near-domain app-family transfer. The
definition is intentionally narrower than open-world generalization,
because the present argument does not require fully unrestricted
transfer in order to be meaningful.

\subsection{Common Confusions to
Eliminate}\label{33-common-confusions-to-eliminate}

Four confusions repeatedly blur the literature and need to be eliminated
before the taxonomy is applied.

\textbf{Confusion 1: workflow cache != procedural memory.} A cached
workflow or approved plan may be highly useful, but it is not
automatically a procedural memory unit in the sense used here. A
workflow cache typically stores a successful route through a task.
Procedural memory, by contrast, should capture the transferable local
rule inside that route, together with the conditions under which it
applies and fails. This distinction is what separates systems such as
MobileGPT, MAGNET, and IntentCUA from the stronger target posed by
reusable procedural memory.

\textbf{Confusion 2: context compression != cross-session episodic
memory.} Summarizing the current trajectory can stabilize long-horizon
execution, but that does not imply the agent now has an episodic memory
architecture. In-task compression is a working-memory function. Episodic
memory requires persistence, retrieval across sessions or tasks, and
some principle for selecting or consolidating prior episodes. Treating
these two as equivalent risks overstating what current long-horizon
systems actually achieve.

\textbf{Confusion 3: stronger training pipeline != deployment-stage
write-back.} Synthetic trajectory generation, post-training, and
training-time self-evolution can all improve GUI agents, but they do not
by themselves create a runtime loop in which failure revises long-term
memory. This distinction is central to the failure-driven write-back
problem. If deployment-stage memory never changes after a failure, then
the system may be self-improving in a broad engineering sense while
still lacking failure-driven write-back in the specific sense that
matters here.

\textbf{Confusion 4: user profile concept != validated user-centric
memory mechanism.} A paper may acknowledge that user history or
preference information matters without providing a mature user-centric
memory loop. This is the current status of much of the personalization
line. Recognizing the object is not the same as showing how it should be
stored, updated, retrieved, and evaluated over time. The distinction
matters because user-centric memory is a confirmed direction, but still
not a solved mechanism.

\begin{center}\rule{0.5\linewidth}{0.5pt}\end{center}

\section{Taxonomy}\label{4-taxonomy}

\subsection{Dimension I: What Is
Learned}\label{41-dimension-i-what-is-learned}

The first dimension asks what kind of experience is being accumulated.
This question matters because not every useful artifact should be
treated as the same memory object. In the current taxonomy, four
categories are needed.

\textbf{Skills} refer to reusable local action policies,
exception-handling patterns, and task-relevant procedures. This is the
category most closely aligned with the present argument, because it is
where experience-dependent procedural knowledge should appear if the
field is truly learning from interaction. A typical GUI example would be
a rule that, in a certain settings layout, the relevant control is often
reached by opening an overflow menu before the visible option list
becomes useful.

\textbf{World} refers to relatively stable environment knowledge such as
app structure, UI semantics, and interface regularities. This category
includes knowledge such as which navigation tabs exist, which screen
regions are functionally stable, or how a site organizes states and
transitions. World knowledge is useful, but it is not automatically
procedural. It helps the agent orient itself; it does not by itself
specify which action policy should be taken under failure or partial
observability.

\textbf{Episodes} refer to concrete past task histories, including
successful and failed trajectories. This category is important because
many systems begin by storing or compressing episodes before they can
abstract rules from them. But episodic storage should be treated as a
substrate rather than as the final answer to procedural reuse. A past
trajectory may be worth retrieving without yet being the right memory
unit.

\textbf{User} refers to user-specific preferences, habits, and
personalized operating tendencies. This category matters because some
GUI tasks are underdetermined without user history. Yet it remains
peripheral to the present argument because the best available evidence
still concerns access to preference history more than mature user-memory
maintenance.

This first dimension therefore separates four different reasons to store
experience. The distinction is not merely descriptive. It prevents the
survey from labeling every external artifact as "procedural memory" and
clarifies why the present argument is specifically centered on
\texttt{Skills}, with \texttt{Episodes} as supporting substrate,
\texttt{World} as structural context, and \texttt{User} as a separate
personalization-oriented direction.

\subsection{Dimension II: When Is Experience Written
Back}\label{42-dimension-ii-when-is-experience-written-back}

The second dimension asks when experience is written back or
consolidated. Four stages are needed here as well: \texttt{pre-task},
\texttt{in-task}, \texttt{post-task}, and \texttt{cross-task}.

\texttt{Pre-task} learning refers to exploration or preparation before
the target task is executed. App exploration and app-specific knowledge
construction belong here. This stage is already well represented in GUI
agents and explains why systems such as AppAgent and MobileGPT can
outperform purely stateless prompting. But pre-task learning often
remains app-bound and does not guarantee that later task outcomes will
revise the memory itself.

\texttt{In-task} learning refers to real-time reflection, compression,
or correction while the task is unfolding. This includes working-memory
updates, step-level verification, and local replanning. The field is
increasingly competent at this stage, especially in long-horizon agents.
Still, in-task adjustment is not yet the same as durable experience
accumulation because the updated state often disappears when the episode
ends.

\texttt{Post-task} learning refers to the period after a task outcome is
known, when successful or failed trajectories can be summarized,
abstracted, filtered, or written back. This stage is crucial because it
is the first point at which the agent can compare expected and actual
outcome and decide whether a memory artifact should be added or revised.
MAGNET and related systems provide the clearest GUI-side early examples
here, but the literature still rarely reaches rule-level generalization
or failure-driven revision.

\texttt{Cross-task} learning refers to the persistent maintenance and
reuse of experience across later tasks. This is where the present
argument becomes most demanding. The central bottleneck is not whether
any learning occurs at all, but whether post-task evidence can be
consolidated into cross-task reusable memory. This is the exact bridge
the field still lacks: many systems can explore, reflect, or summarize,
but far fewer can preserve the resulting knowledge in a form that
remains retrievable, revisable, and useful later.

For this reason, the survey\textquotesingle s main stance is not that
pre-task or in-task learning are unimportant. It is that the decisive
missing link lies between \texttt{post-task} and \texttt{cross-task}. A
system that cannot cross that boundary may still be adaptive in the
short term, but it does not yet solve the experience-reuse problem
targeted in this survey.

\subsection{Dimension III: How Is Experience
Reused}\label{43-dimension-iii-how-is-experience-reused}

\subsubsection{Match Operator}\label{431-match-operator}

The third dimension asks how experience is reused, and here the taxonomy
must be explicitly split into two layers. The first layer is the
\textbf{Match Operator}, which describes how the system finds relevant
prior experience.

\texttt{Exact\ match} refers to direct reuse when the current task,
screen pattern, or sub-task closely matches a previously stored case.
This is common in task recall and app-internal reuse settings. It is
often the easiest form of memory reuse to validate, but it also has the
weakest transfer potential.

\texttt{Similarity\ retrieval} refers to retrieving prior experience
based on approximate semantic, structural, or multimodal similarity.
This is more flexible than exact match and already appears in workflow
and insight retrieval systems. However, it introduces the classic
problem of interference: retrieved memories may look relevant while
encoding the wrong abstraction or the wrong scope.

\texttt{Rule\ generalization} refers to reusing an abstract procedural
rule that is no longer tied to one exact past instance. This is the most
important operator in the present framing because it is what separates
reusable procedural memory from replay or nearest-neighbor recall. If
the field cannot support rule-level generalization, then cross-task
reuse will remain largely equivalent to better cache lookup.

\subsubsection{Application Carrier}\label{432-application-carrier}

The second layer is the \textbf{Application Carrier}, which describes
how the retrieved experience actually influences behavior once it has
been found.

\texttt{Context\ injection} means inserting the retrieved artifact into
the planner or actor context. This is the most common carrier in current
systems because it is easy to implement and compatible with prompt-based
agents. But it is also the easiest to overuse, and it often blurs the
difference between passive reminder and operational constraint.

\texttt{External\ tool\ call} means using memory as a callable module,
retriever, or workflow executor rather than as plain prompt context.
This carrier is rare in current GUI memory work but conceptually
important because it creates a clearer separation between long-term
experience and the main reasoning prompt.

\texttt{Policy\ constraint} means using memory to shape action selection
more directly, for example as a verifier prior, reranking signal, or
decision constraint. This carrier is especially relevant when the goal
is not merely to remind the agent of prior experience, but to make that
experience consistently affect execution.

\texttt{Memory\ rewrite} means that experience is not only applied to
the current task but also used to modify the memory store itself. This
is the carrier most relevant to failure-driven write-back. It captures
the fact that memory maintenance is itself a mechanism of reuse: the
system uses current evidence to decide how the next use of memory should
change.

We separate matching from application because prior taxonomies often
conflate how experience is found with how it is enforced. This
distinction is necessary here because
\texttt{similarity\ retrieval\ ×\ context\ injection} and
\texttt{rule\ generalization\ ×\ policy\ constraint} may both be called
"memory use," yet they represent very different capability levels. One
is closer to reminding the model of a relevant prior case; the other is
closer to letting stored experience directly shape the decision.

\subsection{Interpretive Validity of the
Taxonomy}\label{44-interpretive-validity-of-the-taxonomy}

This taxonomy is intended as an analytical framework rather than a
single-label classification of whole systems, so its usefulness depends
on whether its dimensions remain conceptually distinct at the level of
the capability being analyzed. Under that interpretation, \texttt{What}
and \texttt{When} remain sufficiently separate to identify both the kind
of experience being stored and the stage at which it is consolidated.
The same applies to \texttt{How} once matching is separated from
application: distinguishing \texttt{Match\ Operator} from
\texttt{Application\ Carrier} avoids conflating how experience is found
with how it influences behavior, and makes the framework stable enough
for comparative analysis.

The main caveat is that system-level mapping is not single-label. One
GUI agent may contain multiple memory units, and those units may occupy
different cells of the taxonomy. This does not invalidate the taxonomy.
It simply means the right object of analysis is the capability pattern
rather than the whole paper taken as one atomic category. Under that
interpretation, the framework is both mutually discriminative enough to
be useful and collectively complete enough for the present argument.

\subsection{Coverage Pattern of Current
Systems}\label{45-coverage-pattern-of-current-systems}

When current systems are mapped into this taxonomy, three conclusions
emerge immediately. First,
\texttt{post-task\ /\ Skills\ /\ similarity\ retrieval\ ×\ context\ injection}
is no longer empty. It is already represented by systems such as MAGNET
~\parencite{sun2026magnet}, which shows that GUI agents can indeed summarize
successful experience after a task and reuse it later. Second,
\texttt{post-task\ /\ Skills\ /\ rule\ generalization} remains sparse.
Current systems still struggle to turn task-completion experience into
compact procedural rules that generalize beyond the original workflow
instance. Third, \texttt{cross-task\ /\ Skills} and especially the
\texttt{memory\ rewrite} carrier remain the most structurally
underdeveloped parts of the space.

This coverage pattern is exactly what justifies the
survey\textquotesingle s central argument. The field no longer needs a
generic argument that memory is useful. It needs a sharper argument that
the dominant current solutions stop at relatively weak forms of reuse:
they retrieve similar successful experience, often inject it into
context, and sometimes maintain workflow-level artifacts, but they
rarely generalize those artifacts into revisable cross-task rules. The
taxonomy therefore does more than organize the literature. It identifies
where the literature is already credible, where current evidence remains
limited, and where the most defensible next contribution lies.

\begin{center}\rule{0.5\linewidth}{0.5pt}\end{center}

\section{Related Work Through the
Taxonomy}\label{5-related-work-through-the-taxonomy}

Rather than reviewing GUI memory systems as a flat list, we organize
prior work by the role memory plays in the agent loop. This matters
because different systems externalize very different kinds of
experience. Some help the agent become familiar with an app before
acting; some stabilize the current trajectory; some store workflows,
plan templates, or structural graphs; and only a very small subset
begins to approach post-task consolidation. Seen through this taxonomy,
the literature already contains several strong partial solutions, but
they correspond to different capability types and should not be
collapsed into a single notion of "agent memory."

\subsection{Pre-Task Exploration and App-Bound
Knowledge}\label{51-pre-task-exploration-and-app-bound-knowledge}

Early GUI memory systems primarily use memory to support environment
familiarization. AppAgent~\parencite{zhang2023appagent} is the clearest
instance of this pattern. It explores an app, compiles discovered
action effects into a per-app knowledge document, and then reuses that
document during deployment. This is already an important departure from
purely stateless operation: the system proves that lightweight
exploration can create useful external memory without fine-tuning.
However, the stored object remains an app-bound semantic artifact. The
document captures which elements tend to do what inside one app, but it
does not yet encode a revisable procedural rule with explicit trigger
conditions, constraints, failure signals, or scope beyond that app.

MobileGPT~\parencite{lee2024mobilegpt} pushes this line further by
turning memory into a hierarchical graph over tasks, sub-tasks, and
actions. This is one of the earliest GUI-native examples where past
experience is clearly reused as procedural structure rather than as flat
notes. Crucially, sub-task knowledge can be shared across tasks within
the same app, which makes MobileGPT a strong precursor to procedural
memory. Yet its memory remains per-app and per-device, and its repair
loop still relies on human intervention rather than autonomous
consolidation. Relative to AppAgent, MobileGPT raises the level of
abstraction from knowledge documents to reusable sub-task schemas;
relative to the present argument, it still stops before
\texttt{post-task\ -\textgreater{}\ cross-task} rule abstraction and
failure-driven write-back. The broader lesson from this group is that
early GUI memory work is not memory-free. The field has already learned
how to externalize app knowledge and app-internal skills. What remains
open is whether these app-bound artifacts can be turned into compact,
revisable, and transferable procedural units.

\subsection{In-Task Reflection, Compression, and Local
Repair}\label{52-in-task-reflection-compression-and-local-repair}

Another major line improves long-horizon execution not by storing
reusable rules, but by making the current run more manageable.
MobileAgent-v2~\parencite{wang2024mobileagentv2} is representative
here. Its memory unit is explicitly scoped to the current task and
works together with planning and reflection agents to preserve relevant
context, track progress, and classify errors. This is highly effective
for multi-step execution because it reduces drift and supports local
recovery, but the memory disappears with the task. The paper is
therefore useful precisely because it clarifies the boundary: in-task
working memory and reflection can substantially improve execution, yet
they do not automatically create a durable experience layer for future
tasks.

This distinction becomes even clearer in later long-horizon augmentation
systems such as GUI-Owl-1.5 /
MobileAgentV3.5~\parencite{xu2026mobileagentv35} and
M2~\parencite{yan2026m2}. Both show that context compression and
retrieval can materially improve practical performance. M2 is
especially informative because it combines recursive trajectory
summarization with external insight retrieval and demonstrates that
high-level experience can help without retraining the backbone. Yet its
external memory is still an insight bank derived mainly from successful
trajectories, not a deployment-generated memory that is continuously
revised after failure. These methods therefore provide strong evidence
that experience representation matters, but they mainly solve context
management and retrieval-time augmentation rather than the full
write-back lifecycle. Once temporary compression is separated from
persistent procedural accumulation, it becomes easier to see why
current long-horizon gains still leave the main research problem
unresolved.

\subsection{Post-Task Workflow Memory as the Current Best GUI-Side
Partial
Solution}\label{53-post-task-workflow-memory-as-the-current-best-gui-side-partial-solution}

The strongest GUI-side partial solutions emerge when experience is
lifted from local interaction history into workflow, structure, or
intent-level memory. MAGNET~\parencite{sun2026magnet} is currently the
clearest representative of this transition. By introducing stationary
memory for stable functional UI semantics and procedural memory for
abstract workflows, it shows that GUI agents can benefit from memory
objects richer than app notes or transient summaries. More importantly,
it demonstrates that memory can help cope with UI drift and app-version
changes, which directly targets a central weakness of app-bound
prompt-only agents. However, MAGNET still relies mainly on successful
trajectories for memory construction, and its procedural memory remains
largely workflow-level and per-app in practice. It therefore represents
an important step in the taxonomy, but still only a limited one
relative to the present argument.

ActionEngine~\parencite{zhong2026actionengine} and
IntentCUA~\parencite{lee2026intentcua} extend this shift in two
different directions. ActionEngine turns memory into a state-machine
graph that supports one-shot program synthesis and deterministic
execution. This is a strong example of structural memory for planning,
but the stored object is still primarily topology and executable
template rather than revisable experience-delta rules. IntentCUA
instead abstracts traces into intent groups, subgroup skills, and
shared plan memory. Its contribution is to show that GUI memory can
operate above the atomic-action level, which is conceptually close to
procedural abstraction. Yet the resulting memory is still closer to
approved plan cache and skill hints than to an open-ended long-term
memory that can be rewritten after failure. Taken together, these
systems confirm that the field has already started moving toward more
structured and reusable external memory, but that the dominant forms
remain workflow-level, app-bound, or success-biased. The remaining gap
is therefore no longer whether memory helps, but whether it can support
finer-grained rule abstraction and principled revision.

\subsection{Transferable Precedents from Memory and Self-Evolving
Agents}\label{54-transferable-precedents-from-memory-and-self-evolving-agents}

Many of the strongest design precedents for the missing GUI capability
currently come from adjacent fields rather than from GUI-native systems.
On the memory side, Generative
Agents~\parencite{park2023generativeagents} and
MemGPT~\parencite{packer2023memgpt} establish two complementary
lessons: raw experience must be selectively retrieved and often
reflected upon, and long-term memory requires explicit management
rather than blind context accumulation.
A-MEM~\parencite{xu2025amem} extends this argument by showing that
memory need not be append-only; later evidence can revise earlier
memory objects. These systems do not yet solve GUI procedural reuse
because their substrates are mostly textual and their actions are not
grounded in dynamic interfaces. Still, they remove several conceptual
objections. They show that persistent retrieval,
reflection-triggered abstraction, self-directed memory operations, and
even memory rewrite are all technically viable.

On the self-evolving side, AWM~\parencite{wang2024awm},
ExpeL~\parencite{zhao2024expel},
SkillRL~\parencite{xia2026skillrl}, and
AgentKB~\parencite{tang2025agentkb} provide even stronger precedents
for the present argument. AWM demonstrates that historical trajectories
can be abstracted into reusable workflows rather than replayed
verbatim. ExpeL shows how success-failure contrasts can maintain an
editable insight pool. SkillRL
goes one step further by extracting generalizable skills from successes
and failures and co-evolving them with policy. AgentKB shows that
transferable experience can improve heterogeneous agents when retrieval
is filtered and disagreement is controlled. None of these papers
directly answers the GUI problem, because they are validated mainly in
textual or web settings and do not address grounded visual interaction,
interface drift, or app-family transfer. But together they make the main
conclusion difficult to avoid: the strongest ingredients for reusable
experience already exist, just not yet in a GUI-grounded form.

\subsection{User-Centric Personalization as a Peripheral but
Confirmed
Line}\label{55-user-centric-personalization-as-a-peripheral-but-confirmed-line}

User-centric memory remains peripheral to the present argument, but it
is no longer a speculative direction. OS-Copilot /
FRIDAY~\parencite{wu2024oscopilot} already make room for user profile
information inside a broader memory architecture, which is important
because it shows that system designers recognize user history as a
distinct memory object rather than merely another prompt field. At the
same time, these systems do not yet provide a cleanly validated
user-centric write-back loop. Their procedural memory is primarily
tool-centric, and their evidence for personalization remains indirect.

Persona2Web~\parencite{kim2026persona2web} makes the missing piece explicit by
showing that current web agents fail catastrophically on ambiguous
personalized tasks when user history is absent, and still perform poorly
even when history is accessible. The benchmark therefore sharpens the
problem: having a history store is not the same as using it effectively,
and using it effectively is not the same as maintaining it over time.
For this survey, user-centric memory should be treated as a confirmed
secondary direction with important evaluation implications, but not as
the most feasible or best-evidenced centerpiece of the present research
agenda.

\begin{center}\rule{0.5\linewidth}{0.5pt}\end{center}

\section{Gap Analysis}\label{6-gap-analysis}

\subsection{Main Gap: Procedural Memory for
GUI}\label{61-main-gap-procedural-memory-for-gui}

Current GUI agents still lack a reusable procedural memory that can
transform post-task experience into cross-task applicable rules. This
claim does not mean the field lacks memory altogether. The progression
from AppAgent~\parencite{zhang2023appagent} to
MobileGPT~\parencite{lee2024mobilegpt} to
MAGNET~\parencite{sun2026magnet},
ActionEngine~\parencite{zhong2026actionengine}, and
IntentCUA~\parencite{lee2026intentcua} shows that GUI systems have
already learned to externalize app knowledge, sub-task structure,
workflows, intent abstractions, and structural state graphs. The
problem is that these artifacts are usually tied to a specific app,
site, workflow family, or cached plan regime. Their memory objects
remain either too coarse, too local, or too rigid to serve as general
\texttt{experience-delta} rules that can be retrieved and reused across
future tasks with clear triggers, constraints, and scope boundaries.

This is why the current best partial solutions do not yet close the gap.
MobileGPT~\parencite{lee2024mobilegpt} demonstrates app-internal
procedural reuse, but its memory is explicitly per-app and per-device.
MAGNET~\parencite{sun2026magnet} demonstrates workflow-level procedural
memory and dynamic evolution, but its memory still depends on
successful trajectories and remains constrained by per-app workflow
abstraction. ActionEngine~\parencite{zhong2026actionengine} proves that
structural memory can dramatically improve planning, yet its graph
stores site topology and executable templates rather than revisable
experiential rules. IntentCUA~\parencite{lee2026intentcua} shows that
skill abstraction can rise above atomic action traces, but its memory
is still closer to user-approved plan cache and subgroup hints than to
continuously consolidated procedural memory. These systems collectively
show that the field has reached the edge of the problem, not its
solution.

The importance of this gap becomes even clearer once adjacent precedents
are considered. AWM~\parencite{wang2024awm} and related self-evolving
systems already demonstrate that trajectories can be distilled into
reusable workflow knowledge in textual or web domains. The unresolved
jump is therefore not conceptual feasibility but GUI grounding: how to
represent procedural knowledge that is abstract enough to transfer
beyond one exact screen flow, yet grounded enough to remain operational
under visual variation, hidden preconditions, and layout drift. This is
precisely why procedural memory for GUI remains the central gap. It sits
at the intersection of what GUI-native systems have almost achieved and
what adjacent fields have already partially validated.

\subsection{Main Gap: Failure-Driven
Write-Back}\label{62-main-gap-failure-driven-write-back}

Current GUI agents may accumulate successful traces or benefit from
stronger training flywheels, yet they still lack a deployment-stage
mechanism that rewrites memory based on failure. This is a distinct
problem from merely having memory. AppAgent~\parencite{zhang2023appagent}
and MobileGPT~\parencite{lee2024mobilegpt} demonstrate memory
construction,
but neither provides autonomous failure-driven revision.
MobileAgent-v2~\parencite{wang2024mobileagentv2} is especially
revealing because it explicitly classifies erroneous and ineffective
operations during execution, yet this signal is not consolidated into
persistent memory. MAGNET~\parencite{sun2026magnet} goes further by
evolving memory over time, but still requires successful trajectories
for memory construction and therefore leaves failure as mostly a
filtering condition rather than a source of learnable correction.

The gap is further sharpened by recent training-oriented self-evolution
work. Systems such as GUI-Owl-1.5 /
MobileAgentV3.5~\parencite{xu2026mobileagentv35} and related
training-side flywheels show that GUI agents can improve through
synthetic data, trajectory refinement, and stronger post-training
pipelines. But these improvements primarily happen before deployment or
outside the runtime memory loop. They do not answer whether a deployed
agent can observe that a previously stored rule no longer works,
diagnose why it failed, and then edit, split, downweight, or rewrite
that rule for future tasks. Training-time evolution and deployment-time
write-back should therefore be treated as related but non-equivalent
capabilities.

Adjacent fields again provide the strongest precedents without yet
closing the GUI gap. AWM~\parencite{wang2024awm} and
ExpeL~\parencite{zhao2024expel} show that success-failure contrasts can
yield reusable experience abstractions. SkillRL~\parencite{xia2026skillrl}
shows that failures can be distilled into generalized corrective skills
rather than stored as raw cautionary traces.
A-MEM~\parencite{xu2025amem} shows that memory can be revised rather
than merely appended. Yet none of these systems addresses the grounded
GUI problem where failures arise from visual ambiguity, hidden
interface state, transient pop-ups, or app-specific control logic. This
is why failure-driven write-back remains a central research problem
rather than an implementation detail of procedural memory. Without
failure-driven write-back, even a good procedural memory will tend to
degenerate into a success-case cache.

\subsection{Supporting Gap: Episodic Memory for
GUI}\label{63-supporting-gap-episodic-memory-for-gui}

Episodic memory is best treated as supporting infrastructure for
retrieval and evaluation, rather than the central object of the present
argument. Generative Agents~\parencite{park2023generativeagents} and
MemGPT~\parencite{packer2023memgpt} already show why long-term storage,
selective retrieval, and memory management matter, while GUI-side
results such as
M2~\parencite{yan2026m2} suggest that summarized history and insight
retrieval can materially improve long-horizon performance. But episodic
memory alone does not automatically produce reusable procedural
knowledge. Raw history helps preserve context; it does not by itself
explain which local strategies should transfer, when a previously
successful tactic has become invalid, or how a failure should modify an
existing rule. For this reason, episodic memory is better positioned as
an enabling layer beneath procedural memory and failure-driven
write-back than as the survey\textquotesingle s central research
object.

\subsection{Extension Gap: User-Centric
Memory}\label{64-extension-gap-user-centric-memory}

User-centric memory is a high-value secondary direction, but it
currently sits outside the most feasible and best-evidenced contribution
path. OS-Copilot / FRIDAY~\parencite{wu2024oscopilot} show that user
information can be acknowledged as part of an agent memory architecture,
and
Persona2Web~\parencite{kim2026persona2web} demonstrates that
history-dependent personalization is a real and difficult capability gap
rather than a product nicety. However, the available evidence still
mostly concerns access to history and use of preference signals, not a
mature loop for maintaining, revising, and safely deploying
user-centric GUI memory over time. The direction is therefore
confirmed, but still not strong enough to displace the
procedural-memory focus adopted here.

\subsection{Why These Gaps Matter More Than Other
Blanks}\label{65-why-these-gaps-matter-more-than-other-blanks}

Not every empty cell in the taxonomy deserves equal attention. The
reason procedural memory for GUI and failure-driven write-back matter
more than other blanks is not simply that they are empty, but that they
satisfy four stronger conditions simultaneously. They are technically
plausible because adjacent fields have already validated major
sub-components such as workflow abstraction, editable experience pools,
and memory rewrite. They align directly with the current research
question because they target the missing
\texttt{post-task\ -\textgreater{}\ cross-task} experience layer rather
than peripheral capabilities. They are experimentally tractable because
repeated-task, same-app transfer, and near-domain app-family transfer
already offer realistic validation settings. And they carry higher
explanatory value than many alternative blanks because solving them
would clarify whether GUI agents need more model capacity, better
prompts, or a genuinely new memory mechanism.

By contrast, some other blanks are either too weakly evidenced, too far
from current feasibility, or too indirect relative to the central
argument. This is why the survey should not present the taxonomy as an
invitation to fill every vacancy. Its value is diagnostic: it helps
separate structural bottlenecks from merely unoccupied possibilities. On
the current evidence base, reusable procedural memory and failure-driven
write-back remain the two most defensible targets.

\begin{center}\rule{0.5\linewidth}{0.5pt}\end{center}

\section{Methodological
Implications}\label{7-methodological-implications}

\subsection{What Should Be Stored}\label{71-what-should-be-stored}

The literature reviewed so far suggests that the target memory object
should not be a raw trajectory and should not simply restate generic GUI
priors. Raw trajectories are too verbose, too noisy, and too brittle:
they preserve local detail but do not directly explain what should
transfer. Generic priors, on the other hand, do not justify an external
memory mechanism at all, because the model should already know many of
them. What should be stored instead is a compact procedural rule that
captures the part of the experience that was revealed only through
interaction and is likely to matter again.

At minimum, such a memory unit should encode a trigger, a recommended
procedure, the relevant constraints or preconditions, a failure signal
indicating when the stored strategy has become unreliable, and a scope
statement describing where the rule is expected to transfer. This is the
smallest representation that remains faithful to the current gap
analysis. It is richer than app notes, more abstract than raw action
traces, and more editable than frozen workflows. It also directly
answers a recurring objection from the related-work section: if the
stored object cannot specify when it should fail or be revised, it will
likely collapse back into either case-based retrieval or static workflow
caching.

\subsection{How Memory Should Be Retrieved and
Applied}\label{72-how-memory-should-be-retrieved-and-applied}

The survey also suggests that retrieval and application should be
treated as separate design decisions rather than a single mechanism. On
the retrieval side, the natural comparison is among raw trajectory
retrieval, coarse workflow memory, and fine-grained delta procedural
memory. These baselines correspond to three different assumptions about
what makes past experience useful: exact recall of prior episodes, reuse
of higher-level workflows, or reuse of compact procedural deltas. A
convincing follow-on method should show not only that memory helps, but
that the finer-grained representation is the right unit of reuse.

On the application side, the literature strongly suggests that context
injection alone is not enough. Many current systems prepend retrieved
memory to the prompt, which is often the simplest starting point, but
this makes it hard to distinguish between memory as actionable guidance
and memory as passive reminder. The taxonomy therefore motivates
evaluating multiple application carriers: prompt-level context
injection, tool-mediated retrieval, policy-level constraint or
reranking, and memory rewrite as a maintenance carrier. This separation
is not cosmetic. It determines whether the system merely sees prior
experience or can actually let that experience shape planning,
execution, and revision in a stable way.

\subsection{Why Failure Must Trigger Memory
Revision}\label{73-why-failure-must-trigger-memory-revision}

Failure-driven revision emerges from the survey not as an optional
enhancement, but as a core methodological requirement. Without a
write-back path, even a good memory store tends to become a success-case
cache. This is already visible in current GUI systems that benefit from
successful workflows or approved plans but have little principled way to
revise those artifacts when the environment changes or when an old
strategy begins to fail. The consequence is not just stagnation. It is
miscalibrated reuse: the system continues to retrieve an experience
artifact whose original conditions no longer hold.

This is why rewrite operations such as \texttt{edit}, \texttt{split},
\texttt{downweight}, or \texttt{rewrite} should be treated as
first-class design choices rather than implementation detail. Each
operation corresponds to a different hypothesis about why the prior
memory failed. An edit assumes the memory was mostly right but
incompletely specified. A split assumes two contexts were previously
conflated. A downweight assumes the memory remains valid but should be
trusted less broadly. A rewrite assumes the previously stored policy is
no longer the right abstraction. A subsequent method need not commit to
the final write-back policy at the outset, but it should inherit from
this review the principle that failure must be able to change memory,
not merely trigger another attempt.

\subsection{Minimum Evaluation
Package}\label{74-minimum-evaluation-package}

The minimum evaluation package follows directly from the central
argument of the survey. At a baseline level, a future method should be
compared against \texttt{no\ memory},
\texttt{raw\ trajectory\ retrieval}, \texttt{success-only\ memory},
\texttt{failure-aware\ memory\ update},
\texttt{coarse\ workflow\ memory}, and
\texttt{fine-grained\ delta\ procedural\ memory}. These comparisons are
necessary because each baseline eliminates a different alternative
explanation. If the full method only beats \texttt{no\ memory}, it may
simply be exploiting any added context. If it fails to beat workflow
memory, then fine-grained procedural abstraction may not be justified.
If it fails to beat success-only memory, then the write-back argument
remains unproven.

Evaluation should also be staged across three transfer boundaries:
repeated-task reuse, same-app cross-task reuse, and near-domain
app-family transfer. This order matters. Repeated-task improvement tests
whether the system can learn from experience at all. Same-app transfer
tests whether the stored memory generalizes beyond literal task replay.
Near-domain app-family transfer tests whether the learned unit is more
than an app-specific cache. Open cross-app transfer may still be worth
reporting, but the survey suggests it should be treated as a bonus
rather than a precondition for the main contribution. The main empirical
question is not whether memory solves all GUI generalization, but
whether it produces measurable experience gain beyond what the base
model already knows.

\begin{center}\rule{0.5\linewidth}{0.5pt}\end{center}

\section{Project Objectives and Proposed
Methodology}\label{8-project-objectives-and-proposed-methodology}

\subsection{Project Objective}\label{81-project-objective}

The project takes the research question stated in Section~2.3 as its
starting point and turns it into a concrete methodological objective: to
build a GUI-agent memory layer that stores reusable local procedural
rules, retrieves them beyond literal replay, and revises them when
failures reveal missing constraints or incorrect scope.

At the project level, this objective translates into four concrete
commitments. First, the stored memory should capture interaction-derived
procedural knowledge rather than generic GUI priors. Second, the memory
should remain usable across later tasks rather than only within one
trajectory. Third, failures should trigger memory maintenance rather
than merely runtime recovery. Fourth, empirical gains should extend
beyond repeated-task replay and remain visible under more demanding
transfer settings.

\subsection{Proposed Memory Object}\label{82-proposed-memory-object}

The proposed method instantiates these commitments through a local
\texttt{experience-delta\ procedural\ rule}, or EDPM. An EDPM is a
compact rule that captures what the agent learned from one local
interaction and when that lesson should transfer. This is different from
storing a raw trajectory, a frozen workflow cache, or a general semantic
note. In the current design, each memory unit records at least five
fields: a trigger condition, a recommended procedure, relevant
constraints or preconditions, a failure signal, and an explicit transfer
scope.

This design choice gives the project a specific target object for
implementation. Instead of asking in general whether memory helps, the
project asks whether this smaller and more revisable unit is a better
basis for cross-task reuse than raw traces or coarse workflow templates.

\subsection{Proposed System Design}\label{83-proposed-system-design}

The planned system contains four core modules. The first is
\textbf{candidate mining}, which identifies segments of post-task
experience that contain reusable procedural regularities rather than
incidental episode detail. The second is \textbf{rule abstraction},
which converts those segments into compact memory units with explicit
conditions of use. The third is
\textbf{retrieval-application separation}, which distinguishes how
memory is matched from how it influences the agent, allowing retrieval
quality and application policy to be evaluated independently. The fourth
is \textbf{failure-aware write-back}, which updates existing memory when
later evidence shows that a previously stored rule was incomplete,
over-broad, or incorrect.

This design is intentionally narrower than several nearby alternatives.
It does not simply repackage app-bound knowledge documents, workflow
memories, or structural state graphs. It also does not treat
success-only accumulation as sufficient evidence of an evolving
experience layer. Finally, it does not equate training-time
self-improvement with deployment-time memory revision. The project is
therefore positioned as a memory-layer contribution rather than as a
broader scaling or post-training strategy.

\subsection{Evaluation Methodology}\label{84-evaluation-methodology}

The evaluation plan operationalizes the methodology from Section~7 in a
project-specific form. For readability, the comparison conditions are
labeled as a compact baseline set rather than repeated in full prose:
\texttt{C0} no memory, \texttt{C1} raw retrieval,
\texttt{C2} coarse workflow memory, \texttt{C3} EDPM with success-only
accumulation, \texttt{C4} EDPM with failure-aware write-back, and
\texttt{C5} EDPM combined with a stronger policy constraint. These
baselines are meant to distinguish genuine procedural reuse from generic
context gain, workflow-level reuse, and the possibility that
failure-aware revision adds no measurable value.

Evaluation will be staged across three transfer levels: \texttt{L1}
repeated-task reuse, \texttt{L2} same-site or same-app cross-task reuse,
and \texttt{L3} near-domain site-family or app-family transfer. The most
important project-level hypotheses are that \texttt{C3} should
outperform both \texttt{C1} and \texttt{C2} on \texttt{L2}, and that
\texttt{C4} should outperform \texttt{C3} on failure-first episodes. A
further guardrail is that the method should preserve positive gain at
\texttt{L2} without producing negative overall transfer at
\texttt{L3}.

\subsection{Initial Benchmark and Expected
Outcome}\label{85-initial-benchmark-and-expected-outcome}

The most appropriate benchmark for the first validation stage is
WebArena, because it provides grounded multi-step web tasks while
allowing transfer to be instantiated as repeated-task reuse, same-site
cross-task reuse, and near-domain site-family transfer. These settings
are narrower than unrestricted cross-app generalization, but they are
methodologically suitable for determining whether the proposed memory
representation captures reusable procedural knowledge rather than one-off
workflow hints.

The expected outcome of the project is not merely a performance increase
on a single benchmark. The intended contribution is a concrete
demonstration that GUI agents can store interaction-derived procedural
rules as revisable memory units, retrieve them beyond literal replay,
and improve them when failures reveal incorrect scope or missing
constraints. If successful, the project should contribute both an
implementable memory design for GUI agents and an evaluation setup for
distinguishing genuine memory gain from prompt expansion, plan caching,
or success-case accumulation.

\begin{center}\rule{0.5\linewidth}{0.5pt}\end{center}

\section{Discussion}\label{9-discussion}

\subsection{What Would Count as Real Memory
Gain}\label{91-what-would-count-as-real-memory-gain}

The central evaluation risk for this line of work is that apparent
memory gains may be misattributed. A system may improve because it has
more context, because it cached a repeated task template, or because the
backbone was already strong enough and simply benefited from better
prompt organization. For this reason, real memory gain in the sense used
by this survey should satisfy three stronger conditions. First, the gain
should exceed what can plausibly be explained by generic base-model
prior alone. The strongest evidence here is not absolute success rate
but selective improvement on cases where the model initially fails or
behaves unstably and later improves after accumulating experience.
Second, the gain should be more than literal cache reuse. If a method
succeeds only when the future task nearly duplicates a stored workflow,
then it has not yet shown that it learned a transferable procedural
rule. Third, the gain should remain revisable under failure. If a memory
artifact can only help when previously successful but cannot be
corrected when wrong, then the system is still closer to a success-case
repository than to an evolving experience layer.

These criteria imply that the strongest future evidence will likely come
from controlled comparisons rather than from one headline benchmark
number. A persuasive system should show repeated-task improvement, some
degree of same-app cross-task transfer, and a measurable advantage for
failure-aware update over success-only accumulation. In the context of
this survey, that is what would justify the claim that memory is
functioning as reusable procedural knowledge rather than as extra prompt
budget.

\subsection{Limits of the Current
Survey}\label{92-limits-of-the-current-survey}

This survey also has several deliberate limits. First, the evidence base
is uneven across sub-areas. Procedural memory for GUI and failure-driven
write-back are comparatively well supported because they can draw on
both GUI-native systems and adjacent precedents from memory and
self-evolving agents. User-centric memory, by contrast, remains
evidence-limited and should still be discussed cautiously. Second,
several peripheral or recently emerging systems remain only partially
integrated into the current literature synthesis. Their absence does not
overturn the core argument, but it does mean that the taxonomy should be
read as a focused synthesis rather than an exhaustive census of the
field.

Third, the survey intentionally does not require fully open cross-app
generalization as a precondition for significance. This is a
methodological choice rather than a claim that open-world transfer is
unimportant. The current evidence suggests that repeated-task reuse,
same-app transfer, and near-domain app-family transfer are already
meaningful and experimentally tractable boundaries. Requiring
unrestricted cross-app transfer too early would risk turning a concrete
and testable problem into a vague benchmark ideal. Finally, the survey
remains stronger as a taxonomy-driven synthesis than as a comprehensive
field census.

\subsection{Outlook}\label{93-outlook}

The main practical value of this survey is that it turns a broad topic
area into a narrower research program centered on procedural memory and
failure-driven write-back. Its next step is not to reopen the problem
framing, but to translate this argument into a concrete method and
evaluation package.

\begin{center}\rule{0.5\linewidth}{0.5pt}\end{center}

\section{Conclusion}\label{10-conclusion}

This survey argues that the next bottleneck for LLM-based GUI agents is
not simply stronger perception, larger models, or longer context, but a
missing experience layer: reusable, revisable procedural knowledge
acquired through interaction. By organizing the field through
\texttt{What} is learned, \texttt{When} experience is written back, and
\texttt{How} it is matched and applied, we show that current systems
have already moved beyond purely stateless execution, but have done so
unevenly. App-bound knowledge documents, working-memory compression,
workflow memory, structural state graphs, and intent abstractions all
provide useful partial solutions. Yet these solutions remain
concentrated in the better-developed but still limited parts of the
taxonomy and still rarely deliver rule-level generalization, persistent
cross-task reuse, or principled memory rewrite after failure.

This diagnosis leads to a more focused research agenda than the broad
claim that "GUI agents need memory." The most defensible research focus
is the combination of procedural memory for GUI and failure-driven
write-back. In this framing, episodic memory serves mainly as supporting
substrate and evaluation infrastructure, while user-centric memory
remains an important but still secondary direction. The resulting shift
is significant because it turns the problem from an abstract capability
wish into a more testable target: can a GUI agent store
interaction-derived procedural rules, retrieve them beyond literal
replay, and revise them when they fail?

The practical value of this survey is therefore twofold. As a literature
contribution, it provides a taxonomy and gap analysis that make the
current field easier to read without collapsing distinct memory
functions into a single category. As a research transition document, it
narrows subsequent method design toward a concrete implementation and
evaluation package: structured procedural memory units, explicit
separation between retrieval and application, and write-back mechanisms
that treat failure as a source of memory maintenance rather than only as
a runtime error. If this framing is correct, the next substantive
advance in GUI-agent research will come less from scaling model capacity
alone and more from building a grounded experience layer that can
persist, transfer, and be corrected over time.

\begin{center}\rule{0.5\linewidth}{0.5pt}\end{center}

\section{References}

\nocite{*}
\printbibliography[heading=none]

\end{document}
